% ID: SEL-D0007-BB-E04
% ESTADO: cotejada. Segunda pasada visual completa y borrador compilado revisados.
% Fuente: Selectividad/matematicas_aplicadas/ORD_TITULAR_MATEMÁTICAS APLICADAS CCSS-JUNIO-IMPRESO.pdf
% SHA-256: c42a605501a277fe7194b964ad92ad66047be9a62dc254fa2219d756d3d71317
% Página: 1; bloque B; ejercicio original 4.

\begin{enumerate}[label=\alph*)]
  \item \textbf{(1 punto)} Calcule la derivada de las siguientes funciones:
  \[
    f(x)=\ln\left(\frac{x-1}{x+1}\right)
    \qquad g(x)=x^3\cdot e^{2x^2}
  \]
  \item \textbf{(0.7 puntos)} Represente gráficamente la parábola
  $h(x)=x^2+x+1$, indicando el vértice y los puntos de corte con los ejes
  coordenados.
  \item \textbf{(0.8 puntos)} Calcule el área del recinto limitado por la
  gráfica de $h(x)=x^2+x+1$, el eje de abscisas y las rectas
  $x=-\dfrac{1}{2}$ y $x=0$.
\end{enumerate}
